% ---------------------------------------------------------------------------
% The D_4 thread, part IV: spectroscopy.
% ---------------------------------------------------------------------------
\section*{Exercises: the \texorpdfstring{$D_4$}{D4} lattice, IV}

Everything in this chapter --- interpolating operators, momentum projection,
variational methods --- carries over to $D_4$, but the two-sublattice structure
of Exercise~\ref{ex:d4-geom}(d) leaves a visible fingerprint on the momentum
labels. The following exercise works it out.

\problemdraft[the low-momentum shell comparison in (d) was checked numerically for $L=8$ only; confirm the general statement before relying on it]{\label{ex:d4-momproj}\textbf{Spatial momentum projection on $D_4$.}
Momentum projection is the usual sum over a time slice,
\begin{equation}
  \widetilde{O}(t,\mathbf{p}) \;=\;
  \sum_{\mathbf{x}\,\in\,\text{slice }t} e^{-i\mathbf{p}\cdot\mathbf{x}}\,O(\mathbf{x},t),
\end{equation}
taken over the actual site coordinates.
\begin{enumerate}
\item[(a)] Show that the slice at fixed $x_4=t$ is a face-centred cubic lattice
  --- the even-sum sublattice of $\mathbb{Z}^3$ for $t$ even, the odd-sum coset
  for $t$ odd --- and that a periodic $L^3$ box with $L$ even carries $L^3/2$
  sites per slice.
\item[(b)] Show that the dual of the slice is the body-centred lattice
  $E^{*}=\mathbb{Z}^3\cup(\mathbb{Z}^3+(\tfrac12,\tfrac12,\tfrac12))$, so that
  momenta are defined modulo $2\pi E^{*}$ and the Brillouin zone is a truncated
  octahedron. Deduce that there are exactly $L^3/2$ distinct momenta, matching
  the site count, and that
  $\mathbf{p}$ and $\mathbf{p}+(\pi,\pi,\pi)$ define the \emph{same} plane wave
  on a slice.
\item[(c)] The same phase is $-1$ on the odd slices. Show that consequently
  \begin{equation}
    C\bigl(t,\mathbf{p}+(\pi,\pi,\pi)\bigr) \;=\; (-1)^t\,C(t,\mathbf{p}),
  \end{equation}
  so the two labels are indistinguishable on any single slice but differ by an
  alternating sign across slices. Relate this to
  Exercise~\ref{ex:d4-transfer}: which transfer matrix is the natural object,
  and for which spatial momenta does the free dispersion relation fail to have a
  positive-energy solution? What does this imply about the fit form for a
  correlator at general $\mathbf{p}$?
\item[(d)] Show that the point group of the slice is the full cubic group
  $O_h$ of order $48$, so that irrep tables, little groups and subduction rules
  are unchanged. Then compare the low-lying momentum shells with the cubic case
  and say where, if anywhere, the halved momentum set makes itself felt.
\item[(e)] How is gauge-covariant smearing built on such a slice?
\end{enumerate}}
{(a) On $D_4$ the coordinates satisfy $\sum_{i}x_i \equiv t \pmod 2$, which is
the stated alternation; each coset of the even-sum sublattice of $\mathbb{Z}^3$
is fcc, with the $12$ spatial links of Exercise~\ref{ex:d4-geom}(d) as nearest
neighbours. Exactly half the points of the $L^3$ box satisfy the parity
condition, giving $L^3/2$; $L$ must be even for the periodic identification to
respect it.

(b) For $\mathbf{v}$ to satisfy $\mathbf{v}\cdot\mathbf{x}\in\mathbb{Z}$ on all
of $E$ it suffices to test the generators: $(1,1,0)$ and permutations give
$v_i+v_j\in\mathbb{Z}$, and $(2,0,0)$ gives $2v_1\in\mathbb{Z}$, so the $v_i$
are either all integers or all half-odd-integers --- the body-centred lattice,
of covolume $\tfrac12$. Counting momenta modulo $2\pi E^{*}$ in a box of side
$L$ gives $L^3/2$, matching the sites. Finally
$e^{i(\pi,\pi,\pi)\cdot\mathbf{x}}=e^{i\pi\sum_i x_i}=+1$ on an even slice, so
$\mathbf{p}$ and $\mathbf{p}+(\pi,\pi,\pi)$ act identically there. (Direct
enumeration at $L=4$ finds $32$ distinct plane waves on $32$ sites, the
momenta pairing up exactly as claimed.)

(c) On an odd slice $\sum_i x_i$ is odd and the same phase is $-1$; since
$\sum_i x_i\equiv t$, the factor is $(-1)^t$ and the relation follows at once.
This is the momentum-space face of the structure found in
Exercise~\ref{ex:d4-transfer}: the transfer matrix maps between the two
sublattices, so $T^2$ is the operator acting on a fixed Hilbert space and $T$
itself may have negative eigenvalues. Those sectors are exactly the spatial
momenta with $\sum_i\cos p_i\le0$, where the free dispersion relation has no
positive-energy solution. Practically, a correlator at general $\mathbf{p}$
should be fitted with both an ordinary decaying term and a $(-1)^t$ oscillating
term,
\begin{equation}
  C(t,\mathbf{p}) \;=\; \sum_n A_n e^{-E_n t} \;+\; (-1)^t \sum_m B_m e^{-E'_m t},
\end{equation}
exactly as one does for staggered fermions --- except that here the oscillation
is geometric in origin and has nothing to do with tastes. Restricting to even
time separations, i.e.\ working with $T^2$, removes it at the cost of halving
the time resolution.

(d) The fcc lattice has the same point symmetry as the simple cubic lattice:
enumerating the integer orthogonal matrices that permute the $12$ nearest
neighbours gives order $48$, the group $O_h$. So the whole apparatus of this
chapter --- irrep tables, the little groups of $\mathbf{p}\neq0$, subduction,
operator construction --- transfers verbatim. What might have been expected to
change does not: enumerating the shells at $L=8$ gives degeneracies
$1,6,12,8,6,24$ for $|\mathbf{p}|^2 = 0,1,2,3,4,5$ in units of $(2\pi/L)^2$, the
same as the cubic case. The reason is that the identification of (b) shifts
momenta by $(\pi,\pi,\pi)$, whose length is of order the zone boundary, so the
folding only affects momenta near the edge of the Brillouin zone; the low-lying
shells used in practice are untouched. Locating the first shell at which the
two lists diverge is a worthwhile check.

(e) Exactly as usual, but with the $12$ spatial links of the slice in place of
the $6$ of a cubic lattice: a Jacobi- or Gaussian-type smearing iterates
$S = S_0 + \kappa\,\Delta_{\text{fcc}}$ with $\Delta_{\text{fcc}}$ the covariant
Laplacian built from those links. Since the slice is fcc, the smearing operator
inherits the better rotational behaviour of that geometry.}
